\documentclass[10pt]{article}
\usepackage[T2A]{fontenc}
\usepackage[utf8]{inputenc}
% \usepackage{paratype}
\usepackage{PTSansCaption}% PT Sans Caption font
\usepackage[english,russian]{babel}
\usepackage{longtable}

\renewcommand*\familydefault{\sfdefault}

\newlength{\mywidth}
\def\symbolWidth#1{\settowidth{\mywidth}{#1}'#1':~\the\mywidth}

\makeatletter
\newcommand{\printfontinfo}{%
  \texttt{\f@encoding/\f@family/\f@series/\f@shape/\f@size pt}%
}
\makeatother

\begin{document}

Нормальный пробел (fontdimen 2): \the\fontdimen 2 \font\\
Растяжение пробела (fontdimen 3): \the\fontdimen 3 \font\\
Сжатие пробела (fontdimen 4): \the\fontdimen 4 \font\\

Параметры шрифта: \printfontinfo

{\bfseries\itshape Жирный курсив: \printfontinfo}

Внешнее имя шрифта:
Текущий шрифт: \texttt{\fontname\font}

{\itshape Курсив: \texttt{\fontname\font}}

Внутреннее имя шрифта в TeX:
\expandafter\texttt\expandafter{\expandafter\string\the\font}

\makeatletter
Размер шрифта normalsize: \f@size pt, интерлиньяж: \f@baselineskip
\makeatother

\tiny

\makeatletter
Размер шрифта tiny: \f@size pt, интерлиньяж: \f@baselineskip
\makeatother

\LARGE

\makeatletter
Размер шрифта LARGE: \f@size pt, интерлиньяж: \f@baselineskip
\makeatother

\vspace{5mm}

\begin{tabular}{|p{30mm}|p{30mm}|}
\selectlanguage{english}{I was going to do something today but I haven't finished doing nothing from yesterday.}
\selectlanguage{russian}{Я собирался кое-что сделать сегодня, но еще не закончил ничегонеделание со вчерашнего дня}
&
% I was going to do something today but I haven't finished doing noth\-ing from yesterday. 
% Я собирался ко\hskip0pt plus.02em\/е-что сделать сегодня, но еще не закончил ничегонеделание со вчерашнего дня
I was go\-ing to do some\-thing to\-day but I haven't fin\-ished do\-ing nothing fr\-om yes\-terday.
Я со\-би\-рал\-ся ко\hskip0pt plus.02em\/е-что сде\-лать се\-го\-дня, но еще не за\-кон\-чил ни\-че\-го\-не\-де\-ла\-ние со вче\-раш\-него дня
\looseness=-3
\end{tabular}

\begin{tabular}[!t]{c|c|c|c|c|c|c|c|c|c}
\hline
\symbolWidth{1} & \symbolWidth{2} & \symbolWidth{3} & \symbolWidth{4} & \symbolWidth{5} &
\symbolWidth{6} & \symbolWidth{7} & \symbolWidth{8} & \symbolWidth{9} & \symbolWidth{0}
\\
\hline
\end{tabular}

\end{document}

